% =====================================================================
\chapter{Plantilla del documento de diseño}
\label{ap:plantilla}
% =====================================================================

La plantilla desarrolla la de M28, §6.2, para el proyecto final, siguiendo la revisión preliminar de diseño del curso 16.842 del MIT~\cite{mit16842} y el planteamiento de sistemas con aprendizaje automático de CS329S~\cite{cs329s,huyen}. Se rellena antes de escribir código y se actualiza en las semanas~2, 5 y~8 (entregable~E1). Cada sección indica qué debe responder, qué evidencia numérica se exige y en qué versión debe estar completa. Las secciones vacías se marcan como «pendiente», no se eliminan.

\section{Ficha resumen}

La primera página del documento es esta ficha, con una o dos líneas por campo.

\begin{center}
\small
\begin{tabularx}{\textwidth}{@{}>{\raggedright\arraybackslash}p{4.6cm}X@{}}
\toprule
Campo & Contenido \\
\midrule
Opción y título & \\
Autor o autores y reparto & \\
Decisión, frecuencia y horizonte & \\
Objetivo (escalar, con unidades) & \\
Restricciones duras & \\
Indicadores que se vigilan & \\
Línea base / alternativa sencilla / techo & tres cifras con unidades y su origen \\
Formulación y resolvedor & \\
Componente aprendido & qué predice, con qué métrica de decisión \\
Salvaguardas y respaldo & \\
Mejora en prueba (IC del 95\,\%) & frente a la línea base y frente a la alternativa sencilla \\
VAN y $\Prob(\VAN<0)$ & \\
Tres riesgos principales (AMFE) & \\
Orden de reproducción y tiempo & \\
\bottomrule
\end{tabularx}
\end{center}

\section{Secciones del documento}

\begin{table}[H]
\centering\footnotesize
\caption{Secciones del documento de diseño, evidencia exigida y versión en que deben estar completas.}\label{tab:plantilla-secciones}
\begin{tabularx}{\textwidth}{@{}>{\raggedright\arraybackslash}p{2.6cm}>{\raggedright\arraybackslash}X>{\raggedright\arraybackslash}p{3.6cm}c@{}}
\toprule
Sección & Preguntas que debe responder & Evidencia numérica & Versión \\
\midrule
1. Contexto y decisión & ¿Qué se decide, cada cuánto, con qué horizonte y con qué información disponible en el instante de decidir? ¿Qué pasa hoy? & frecuencias, horizontes, volumen de decisiones & v0 \\
2. Partes interesadas & ¿Quién decide hoy, quién ejecuta, quién responde de los errores y quién puede parar el sistema? & --- & v0 \\
3. Objetivo y requisitos & Objetivo escalar; indicadores; restricciones duras y blandas; requisitos de latencia, disponibilidad y calidad con márgenes (\cref{tab:plantilla-requisitos}) & tabla de requisitos con umbrales y márgenes & v0 \\
4. Línea base y techo & Coste de la práctica actual, de la mejor alternativa sencilla y con información perfecta; margen máximo alcanzable & tres cifras con su procedimiento & v0 \\
5. Datos y gemelo & Generador (ecuaciones, parámetros, semillas); validación del gemelo; partición y huella de la prueba & pruebas de validación del gemelo & v0 \\
6. Formulación y método & Modelo matemático completo; familia de métodos elegida y alternativas descartadas con su motivo (mapa de M24, §1.3) & tamaño del modelo, constantes big-M justificadas & v0 \\
7. Modelos y validación & Qué se predice; validación temporal; métricas de predicción y de decisión; selección sin sesgo del ganador & métricas en validación & v1 \\
8. Salvaguardas y respaldo & Comprobaciones independientes; política de respaldo; tiempo máximo de cálculo; modos de operación & prueba automática de cada salvaguarda & v1 \\
9. Evaluación fuera de línea & Protocolo con CRN; resultados frente a las tres referencias con intervalos; ablaciones; pruebas de estrés (\cref{tab:plantilla-resultados}) & tabla de resultados & final \\
10. Piloto & Diseño de un piloto real: modo sombra, tamaño (M28, proposición~5.5), duración, criterio de aceptación fijado de antemano & tamaño del piloto & final \\
11. Despliegue y supervisión & API o herramienta, interfaz, registro de decisiones, versiones y vuelta atrás & latencias medidas & final \\
12. Monitorización & Señales, umbrales, acciones y periodo de reentrenamiento (M28, proposición~4.8) & umbrales calibrados & final \\
13. Valor y riesgo & VAN con incertidumbre, CVaR, AMFE (\cref{tab:plantilla-amfe}), encaje normativo (M28, §5.6) & VAN, $\Prob(\VAN<0)$, NPR & final \\
14. Ciclo de vida & Responsable del mantenimiento, revisiones periódicas, criterios de retirada & --- & final \\
\bottomrule
\end{tabularx}
\end{table}

\section{Tablas que se rellenan}

\begin{table}[H]
\centering\footnotesize
\caption{Plantilla de requisitos (sección~3). Se muestra una fila de ejemplo de la opción~(a).}\label{tab:plantilla-requisitos}
\begin{tabularx}{\textwidth}{@{}c>{\raggedright\arraybackslash}X>{\raggedright\arraybackslash}p{2.3cm}>{\raggedright\arraybackslash}p{1.6cm}>{\raggedright\arraybackslash}p{1.6cm}>{\raggedright\arraybackslash}p{2.6cm}@{}}
\toprule
Id. & Requisito & Métrica & Umbral & Margen de diseño & Verificación \\
\midrule
RQ1 & No superar la potencia contratada & potencia horaria máxima & 1000~kW & 950~kW & prueba automática y gemelo \\
\rule{0pt}{14pt}RQ2 & & & & & \\
\rule{0pt}{14pt}RQ3 & & & & & \\
\bottomrule
\end{tabularx}
\end{table}

\begin{table}[H]
\centering\footnotesize
\caption{Plantilla de resultados (sección~9): una fila por política, en el conjunto de prueba, con CRN y al menos 5 semillas.}\label{tab:plantilla-resultados}
\begin{tabularx}{\textwidth}{@{}>{\raggedright\arraybackslash}X cccccc@{}}
\toprule
Política & \makecell{Coste\\medio} & \makecell{Mejora frente\\a línea base} & \makecell{IC 95\,\%\\$t$ pareado} & \makecell{IC 95\,\%\\\emph{bootstrap}} & \makecell{Servicio} & \makecell{Violaciones\\duras} \\
\midrule
\rule{0pt}{13pt}Línea base & & --- & --- & --- & & \\
\rule{0pt}{13pt}Alternativa sencilla & & & & & & \\
\rule{0pt}{13pt}Optimizador con IA & & & & & & \\
\rule{0pt}{13pt}Optimizador sin componente aprendido & & & & & & \\
\rule{0pt}{13pt}Techo (información perfecta) & & & & & & \\
\bottomrule
\end{tabularx}
\end{table}

\begin{table}[H]
\centering\footnotesize
\caption{Plantilla de AMFE (sección~13), según la definición~5.8 de M28.}\label{tab:plantilla-amfe}
\begin{tabularx}{\textwidth}{@{}>{\raggedright\arraybackslash}p{2.6cm}>{\raggedright\arraybackslash}X>{\raggedright\arraybackslash}p{2.2cm}cccc>{\raggedright\arraybackslash}p{2.6cm}@{}}
\toprule
Modo de fallo & Efecto & Causa & $S$ & $O$ & $Dt$ & NPR & Medida \\
\midrule
\rule{0pt}{14pt} & & & & & & & \\
\rule{0pt}{14pt} & & & & & & & \\
\rule{0pt}{14pt} & & & & & & & \\
\bottomrule
\end{tabularx}
\end{table}

\section{Registro de decisiones de diseño}

Cada decisión importante (elegir un resolvedor, descartar un modelo, relajar un requisito) se anota con: fecha; contexto; opciones consideradas; decisión; consecuencias y cómo se comprobarán; y quién la tomó. El registro se adjunta al final del documento y la corrección lo usa para valorar la sección~6 y la defensa.
